\scanpage{272}
\mspage{- 81 -}\hfill IX

opère (de façon non nécessairement fidèle) sur $k'$ (le corps des invariants étant $k$), et sur $Y_{k'}$ da façon compatible avec son action sur $k'$. Or, d'après la théorie de DIEUDONNÉ [16 bis], au $p$-pro-groupe de Barsotti-Tate $Y_{k'}$ est canoniquement associé un module $M$ libre de type fini (de rang égal au rang de $Y_{k'}$) sur l'anneau $W(k')$ des vecteurs de Witt sur $k'$, et $G$ opère évidemment sur $M$ de façon compatible avec ses opérations sur $W(k')$ (via ses opérations sur $k'$). On en déduit une opération $W(k')$-linéaire du sous-groupe d'inertie $\Gamma$ de $G$ sur $M$ (qui est un quotient du groupe d'inertie $I$), d'où une fonction trace sur $\Gamma$ donc sur $I$, qui est la fonction cherchée.Elle est a priori à valeurs dans $W(k')$ (*), mais coïncide avec la fonction analogue définie en termes des $T_\ell$ pour $\ell\ne p$, lorsque l'opération de $G$ sur $Y_{S'}$ s'obtient à partir d'une opération de $G$ sur un schéma en groupes lisse et $p$-divisible $C'$ sur $S'$, en prenant $Y=T_p(C')$ (cf. [25 bis, th.7 p.737]) ; elle est donc dans ce cas à valeurs entières. On procède de façon analogue dans le cas de 4.3 b) pour définir une fonction trace sur $\pi^!$, en utilisant le procédé de la démonstration de 4.3.2 pour "linéariser" l'opération dans $M$ définie par un élément $g\in\pi$ qui se trouve au-dessus d'un $f_q^n$ avec $n\geq0$.

7. \uline{L'extension de Raynaud} $G^{\handwritten{\unsure{\natural|\sharp}}}$ \uline{sur} $S$ \uline{attachée au schéma abélien} $A_K$ \uline{à réduction semi-stable}. \uline{Dualité des schémas abéliens} $B,B'$ \uline{sur} $S$.

7.1. Soit $S$ un schéma noethérien spectre d'un anneau $V$, séparé et complet pour la topologie définie par un idéal\unsure{.} $\underline m$ de $V$. Soit $A$ un schéma en groupes lisse sur $S$, $A^{\unsure{o|0}}$ sa composante neutre, et supposons que $A_o/A_o^o$ soit fini (donc fini étale) sur $S_o$, et que $A_o^o$ soit une extension d'un schéma abélien

\noindent (*) En fait, il est immédiat que cette fonction est toujours à valeurs dans $W(k)$, étant invariante par $\operatorname{Gal}(k'/k)$.

% Heavy reversed arrows at « neutre » and « et que » match the two left arrows on scan 271.
% archive sheet number: 271
% stamp: 399 (lower centre label; earlier 399 partly covered).
